\documentclass[12pt,a4paper]{article}
\usepackage[utf8]{inputenc}
\usepackage[T1]{fontenc}
\usepackage[portuguese]{babel}
\usepackage{amsmath,amssymb,amsthm}
\usepackage{geometry}
\geometry{margin=2.4cm}
\newtheorem{lema}{Lema}
\newtheorem{corolario}{Corolário}
\newtheorem{observacao}{Observação}
\newtheorem{especulacao}{Especulação}

\title{Avaliação do artigo de Sanctuary (arXiv:0908.3219) \\
e exercício Dirac-relativístico no critério de Nyman--Beurling}
\author{}
\date{}

\begin{document}\maketitle

\section{Veredito sobre o artigo de Sanctuary}

O artigo trata de estados não-hermitianos de spin $1/2$, separabilidade
de estados de Bell, e interpretação por ensembles. Não há conexão
matemática com $\zeta$, $\mu$, Nyman--Beurling ou a linha NUMERIA. A
única analogia é estrutural (não-hermitiano $\to$ hermitiano via
ensemble $\leftrightarrow$ esperança Hilbert--Pólya), sem conteúdo
demonstrativo.

\section{Lema distribucional (real)}

\begin{lema}
Para $k\ge 1$, em $\mathcal{D}'(0,\infty)$:
\[
\rho_k'(x) = -\frac{1}{kx^2} + \sum_{n\ge 1}\delta\!\left(x-\frac{1}{kn}\right).
\]
\end{lema}

\begin{corolario}
$\mathcal{M}[\rho_k'](s) = k^{1-s}\zeta(s-1)$ para $1<\Re s<2$, por
continuação analítica em $\mathbb{C}$. Em particular, para
$\mathcal{D}_N(c)=\chi' - \sum_{k\le N}c_k\rho_k'$,
\[
\mathcal{M}[\mathcal{D}_N(c)](s) = -1 - \zeta(s-1)\sum_{k\le N}c_k k^{1-s}.
\]
\end{corolario}

\begin{observacao}
A fórmula injeta $\zeta$ no numerador, mas não fecha o problema: a
convergência de $\sum c_k k^{1-s}$ na linha $\Re s=1/2$ é exatamente
o obstáculo original, disfarçado como critério de Weil.
\end{observacao}

\section{Leitura relativística}

Com $E=1/2$, $p=s-1/2$ e $m^2=\rho(1-\rho)$:
\begin{center}
\begin{tabular}{ll}
$\operatorname{Im}(m_\rho^2)=0$ & $\iff$ $\Re\rho=1/2$ (RH) \\
$|p|<1/2$ (faixa crítica)      & $\iff$ $\Re s\in(0,1)$ \\
$\operatorname{Re}(m_\rho^2)\ge 1/4$ & se RH vale \\
$\operatorname{Im}(m_\rho^2)\ne 0$ & zero "virtual" (massa complexa)
\end{tabular}
\end{center}

\begin{especulacao}
RH $\iff$ todos os zeros têm massa real positiva. Interpretação física:
zeros são "partículas on-shell"; zeros fora da linha crítica são
"virtuais". Nenhum mecanismo novo força a realidade.
\end{especulacao}

\section{Conclusão}

O artigo de Sanctuary não contribui para RH. Os dois lemas acima são
reais e clássicos; a leitura relativística é limpa mas tautológica. A
linha M2.4-C (cálculo com $Q$ maior, trigama, ansatz restrito) é
ordens de magnitude mais promissora, porque é matemática verificável.
Nenhuma prova de RH é apresentada.

\end{document}